\section{A random walk with asymptotically free products}\label{sec:walk}

Throughout this section, $(M,\tau)$ is a finite von Neumann algebra of
type $\mathrm{II}_1$ with separable predual and faithful normal tracial
state.  In particular, $M$ may have a nontrivial center.  We construct a
single symmetric probability measure on $\U(M)$ whose independent long
products approach the joint moments of free Haar unitaries along a
subsequence.  The construction combines finite approximation to freeness
with increasingly rare increments chosen at successive stages.

Recall that unital subalgebras of a tracial von Neumann algebra are
\emph{free} if every alternating product of centered elements from
successively different subalgebras has trace zero.  Here centered means
having trace zero.  A unitary $u$ is \emph{Haar} if
$\tau(u^k)=0$ for every $k\in\Z\setminus\{0\}$.  A free Haar family
$(u_i)$ is equivalently characterized by
\[
 \tau\bigl(w(u_1,\ldots,u_r)\bigr)=0
 \qquad(w\in\mathbb F_r\setminus\{e\}),
\]
for every finite subfamily, where $\mathbb F_r$ denotes the free group on
$r$ specified generators.  The equivalence follows by expanding centered
Laurent polynomials and then taking bounded $2$-norm approximations.
The length $|w|$ is its reduced word length.

For a probability measure $\mu$ on $\U(M)$, let $\mu^{*n}$ be the law
of the ordered product $U_1\cdots U_n$ of independent increments of
law $\mu$.  Symmetry means that $U^{-1}$ has the same law as $U$.

\begin{theorem}\label{thm:walk}
Let $M$ be a finite type $\mathrm{II}_1$ von Neumann algebra with
separable predual and faithful normal tracial state $\tau$.  There
exist a symmetric, countably supported probability measure $\mu$ on
$\U(M)$ and a strictly increasing sequence $n_\ell\to\infty$ with the
following properties.

The positive-mass atoms of $\mu$ contain a countable $2$-norm dense
subset of $\U(M)$, and hence generate $M$ as a von Neumann algebra.
For every $r\in\N$, every nonempty finite set
$\mathcal W\subset\mathbb F_r\setminus\{e\}$, and every $\eta>0$, if
$G_1^{(\ell)},\ldots,G_r^{(\ell)}$ are independent with law
$\mu^{*n_\ell}$, then
\[
 \Pr\left(
   \max_{w\in\mathcal W}
   \left|\tau\bigl(w(G_1^{(\ell)},\ldots,G_r^{(\ell)})\bigr)\right|
   \geq\eta
 \right)\longrightarrow0.
\]
Thus these products have joint free Haar moments in probability.
\end{theorem}

We first obtain finite unitary configurations for which prescribed
words in path products have small traces.  We then place such
configurations at increasingly rare levels of a single measure,
arranging that each finite moment test succeeds with probability
tending to one.

\subsection{Finite requirements from ultrapower freeness}

The freeness input originates in Popa's free-independent sequence
construction~\cite[Lemma~1.4 and Theorem~2.1]{Popa1995}.
We use the following special case of its later ultraproduct
formulation~\cite[Theorem~0.1(a)]{Popa2014}.

For a tracial von Neumann algebra $(P,\tau_P)$ and a free ultrafilter
$\omega$ on $\N$, its tracial ultrapower is
\begin{equation}\label{eq:tracial-ultrapower}
 P^\omega=\ell^\infty(\N,P)/\mathcal I_\omega,
 \qquad
 \mathcal I_\omega
 =\{(x_m):\lim_{m\to\omega}\norm{x_m}_2=0\}.
\end{equation}
It is a finite von Neumann algebra with faithful normal trace
$\tau_P^\omega((x_m)^\omega)=\lim_{m\to\omega}\tau_P(x_m)$.
We identify $P$ with its constant sequences.

\begin{lemma}[Popa's ultrapower freeness]\label{lem:popa-special}
Let $P$ be a $\mathrm{II}_1$ factor and let $B\subset P^\omega$ be a
unital von Neumann subalgebra with separable predual.  There is a diffuse
abelian von Neumann subalgebra of $P^\omega$ free from $B$.
Consequently, for every $d\in\N$, there are Haar unitaries
$u_1,\ldots,u_d\in P^\omega$ such that
$B,W^*(u_1),\ldots,W^*(u_d)$ are free.
\end{lemma}

\begin{proof}
Apply Theorem~0.1(a) of \cite{Popa2014} with every coordinate algebra
and every coordinate subalgebra equal to $P$.  The required
non-intertwining condition holds: a nonzero corner of the diffuse factor
$P$ cannot embed in the scalar relative commutant, or in a finite matrix
amplification of it.  Moreover, $P^\omega$ is a factor, so the relative
commutant of the full ultraproduct subalgebra is $\C1$.  The conclusion
for the separable centered subspace $B\ominus\C1$ is therefore ordinary
scalar-trace freeness. This is the case in which every coordinate
subalgebra equals its ambient factor, discussed after that theorem
and in \cite[Sections~1.5--1.6]{Popa2014}.

A diffuse abelian tracial algebra contains a Haar unitary.  Choose one,
adjoin it to $B$, and repeat.  Finite adjunction preserves separable
predual: a countable generating set gives a separable tracial Hilbert
space by polynomial approximation.  Thus all $d$ choices take place in
the same ultrapower.  At each step the new algebra is free from the
algebra generated at the preceding steps, giving the asserted family.
\end{proof}

We next record why a single fresh label in each path is sufficient.
A \emph{path expression} is a finite ordered product of fixed unitary
coefficients and formal letters $U_j^{\pm1}$.  Occurrences and labels in
such expressions are counted before any simplification.

\begin{lemma}[Distinguished letters]\label{lem:distinguished-letters}
Let $B,W^*(u_1),\ldots,W^*(u_d)$ be free subalgebras of a tracial von
Neumann algebra, with each $u_j$ Haar.  Let
$y_1,\ldots,y_r$ be path expressions evaluated at these unitaries, with
coefficients in $\U(B)$.  Suppose that for each $i$ there is a label
$h_i$ occurring exactly once in the entire list of path expressions,
and that its occurrence belongs to $y_i$.  Then
\[
 \tau\bigl(w(y_1,\ldots,y_r)\bigr)=0
 \qquad(w\in\mathbb F_r\setminus\{e\}).
\]
\end{lemma}

\begin{proof}
The labels $h_1,\ldots,h_r$ are distinct.  Put all nondistinguished
unitaries into the coefficient algebra:
\[
 B_0=W^*\bigl(B,u_j:j\notin\{h_1,\ldots,h_r\}\bigr),
 \qquad
 C=W^*(u_{h_1},\ldots,u_{h_r}).
\]
Then $B_0$ and $C$ are free.  This grouping property follows directly
by expanding centered polynomials in a free family into alternating
centered products; bounded approximation extends it to the generated
von Neumann algebras.  Write
\[
 y_i=A_i u_{h_i}^{\epsilon_i}B_i,
 \qquad A_i,B_i\in\U(B_0),\quad \epsilon_i\in\{-1,1\}.
\]
If $w=x_{i_1}^{\sigma_1}\cdots x_{i_t}^{\sigma_t}$ is nonempty and
reduced, its evaluated product has the form
\[
 c_0 u_{h_{i_1}}^{\epsilon_{i_1}\sigma_1}c_1
 \cdots u_{h_{i_t}}^{\epsilon_{i_t}\sigma_t}c_t,
 \qquad c_0,\ldots,c_t\in\U(B_0).
\]
The displayed distinguished-letter string is itself reduced.  Indeed,
an adjacent inverse pair would require $i_j=i_{j+1}$ and
$\sigma_j=-\sigma_{j+1}$, contrary to the reduction of $w$.

Expand every $c_j$, including the two outside coefficients, as its
scalar part plus its centered part.  In each resulting term, removing
scalar coefficients joins distinguished letters into nonempty
contiguous subwords of that reduced string.  Every such block is a
centered element of $C$.  The remaining factors alternate between
centered elements of $B_0$ and centered elements of $C$, and there is
at least one $C$ factor.  Freeness makes the trace zero.  If there are
no centered coefficient factors, the term is just a scalar multiple of
a nonempty reduced Haar word and again has trace zero.  This proves
the assertion, with no condition on the first and last indices of $w$.
\end{proof}

\begin{lemma}[Simultaneous finite path requirements]\label{lem:finite-path-tests}
Fix a finite set $F\subset\U(M)$, an integer $d\geq1$, and
$\varepsilon>0$.  Consider a finite collection of tests.  Each test
consists of a finite list of path expressions $y_1,\ldots,y_r$ in the
alphabet $F\cup\{U_j^{\pm1}:1\leq j\leq d\}$ and a nonidentity word
$w\in\mathbb F_r$.  Assume that each path in each test has a
distinguished label occurring exactly once in that test's entire list,
as in Lemma~\ref{lem:distinguished-letters}.  The distinguished labels
may differ between tests.  There are $u_1,\ldots,u_d\in\U(M)$ for which
all the tests satisfy
\[
 \left|\tau\bigl(w(y_1(u),\ldots,y_r(u))\bigr)\right|
 <\varepsilon.
\]
\end{lemma}

\begin{proof}
First suppose that $M=P$ is a factor.  The constant coefficient algebra
$W^*(F)\subset P^\omega$ has separable predual.  Apply
Lemma~\ref{lem:popa-special} and then
Lemma~\ref{lem:distinguished-letters}.  The same free $d$-tuple makes
every prescribed trace exactly zero in $P^\omega$, even though different
tests can designate different labels.

Every unitary of $P^\omega$ has unitary representatives.  To see this,
take a bounded self-adjoint logarithm there, lift it to a bounded
sequence of self-adjoint elements of $P$, and exponentiate coordinatewise.
Choose such representatives for the $d$ unitaries.  Each test trace
converges to zero along $\omega$.  Since the family of tests is finite,
the coordinate sets on which their absolute values are less than
$\varepsilon$ have a common member.  Its $d$ coordinate unitaries solve
all the requirements in $P$.

For general $M$, the factor case will give successful tuples in
almost every fiber.  The remaining task is to choose them measurably.
Use the central direct integral
\cite[Sections~III.1.6 and III.5.1]{Blackadar2006}
\[
 M=\int_\Omega^\oplus M_s\,d\nu(s),
 \qquad
 \tau(x)=\int_\Omega\tau_s(x(s))\,d\nu(s),
\]
where $(\Omega,\nu)$ is a probability space and, almost everywhere,
$M_s$ is a $\mathrm{II}_1$ factor with its normalized trace $\tau_s$.
Standard separable direct-integral theory provides countably many
bounded measurable sections $a_j$ whose values generate $M_s$ for
almost every $s$.  After removing one null set, all these properties
and the unitary identities for the finitely many coefficients hold
simultaneously.

There is a countable family $(v_j)$ of measurable unitary sections
which is dense in each fiber unitary group for
$\norm{\cdot}_{2,s}$.  Here is an explicit construction.  Enumerate the
sections
\[
 \exp(iq(a_1,a_2,\ldots)),
\]
where $q$ ranges over self-adjoint $*$-polynomials in finitely many
variables with coefficients in $\mathbb Q+i\mathbb Q$.  This includes
symmetrizations of arbitrary rational $*$-polynomials.  In a fixed
fiber, any unitary is $\exp(ih)$ for a bounded self-adjoint $h$.
Kaplansky density and norm approximation by rational polynomials
approximate $h$ in $2$-norm by the indicated self-adjoint polynomials.
For bounded self-adjoint $a,b$, the integral identity
\[
 e^{ia}-e^{ib}
 =i\int_0^1 e^{i(1-t)a}(a-b)e^{itb}\,dt
\]
implies $\norm{e^{ia}-e^{ib}}_2\leq\norm{a-b}_2$.
Thus these sections have the required fiberwise density.  Each section
is globally bounded, and exponentiation makes its norm one.

Enumerate all ordered $d$-tuples drawn from $(v_j)$ as
$\boldsymbol v^{(1)},\boldsymbol v^{(2)},\ldots$.  Let $E_k$ be the
set of fibers where candidate $k$ satisfies every test with absolute
trace less than $\varepsilon/2$.  These sets are measurable, since
the traces of products of measurable sections are measurable.  They
cover almost every fiber.  Indeed, the factor case gives a successful
tuple in each such fiber, and the requirements are open in the product
$2$-norm topology.  For this last assertion, telescoping products of
unitaries gives
\[
 \norm{z_1\cdots z_L-z'_1\cdots z'_L}_2
 \leq\sum_{j=1}^L\norm{z_j-z'_j}_2,
 \qquad |\tau_s(x)|\leq\norm{x}_{2,s},
\]
and inversion is also a $2$-norm isometry on unitaries.  The countable
dense tuples therefore meet the open set of successful choices.

Take the first successful tuple: on the disjoint measurable sets
\[
 E_k\setminus\bigcup_{j<k}E_j
\]
use $\boldsymbol v^{(k)}$, and use the identity tuple on the null
complement.  The resulting sections $u_1,\ldots,u_d$ are measurable
and unitary, so belong to $M$.  Each test has fiber trace of absolute
value less than $\varepsilon/2$ almost everywhere.  Integrating gives
an absolute global trace at most $\varepsilon/2<\varepsilon$, as
required.
\end{proof}

\subsection{One measure satisfying all moment tests}

\begin{proof}[Proof of Theorem~\ref{thm:walk}]
Choose a $2$-norm dense sequence $(v_j)_{j\geq1}$ in $\U(M)$ and put
\[
 \mu_0=\sum_{j\geq1}2^{-j-1}(\delta_{v_j}+\delta_{v_j^{-1}}).
\]
For $\ell\geq1$, fix the numerical parameters
\begin{align*}
 p_\ell&=2^{-2^\ell},
 &p_0&=1-\sum_{\ell\geq1}p_\ell,\\
 n_\ell&=p_\ell^{-3/2}=2^{3\cdot2^{\ell-1}},
 &K_\ell&=\ell n_\ell,\\
 d_\ell&=K_\ell^4.
\end{align*}
Here $n_\ell$ and $d_\ell$ are positive integers, and $p_0>0$ because
$p_\ell\leq4^{-\ell}$.  At time $n_\ell$ we will sample $\ell$ paths,
choosing level $j$ with probability $p_j$ at each increment.  These
parameters make level $\ell$ appear in every path with high probability
while later levels remain unlikely.  The much larger label count
$d_\ell$ makes repeated level-$\ell$ labels unlikely across the
$K_\ell$ increments.  We construct finite symmetric measures
\[
 \mu_\ell=\frac1{2d_\ell}
   \sum_{j=1}^{d_\ell}(\delta_{u_{\ell,j}}
                            +\delta_{u_{\ell,j}^{-1}})
\]
successively, and finally set
\[
 \mu=p_0\mu_0+\sum_{\ell\geq1}p_\ell\mu_\ell.
\]

At stage $\ell$, choose a finite symmetric set $F_\ell$ of base atoms
with $\mu_0(F_\ell)>1-K_\ell^{-2}$.  The coefficient alphabet consists
of $F_\ell$ and all $u_{j,h}^{\pm1}$ from stages $j<\ell$; it is
finite.  Introduce $d_\ell$ unknown unitary labels.  Consider every
ordered list of $\ell$ paths of length $n_\ell$ in this alphabet and
the signed unknown labels such that:
\begin{enumerate}
 \item each path contains at least one unknown label;
 \item no unknown label occurs twice anywhere in the entire list,
       irrespective of sign.
\end{enumerate}
For each such list and every nonempty reduced word of length at most
$\ell$ on $\ell$ generators, require its evaluated trace to have
absolute value less than $1/\ell$.  There are finitely many lists and
finitely many words: all alphabet sizes, path lengths, and word lengths
have been fixed before making this choice.  Designating any one unknown
occurrence in each path puts every test under
Lemma~\ref{lem:finite-path-tests}.  Apply that lemma to choose
$u_{\ell,1},\ldots,u_{\ell,d_\ell}$ satisfying all these requirements
simultaneously.  This completes the induction.

We estimate probabilities on a labeled sampling space.  An increment
first chooses a level $j$ with probability $p_j$.  At level zero it
chooses an atom with law $\mu_0$; at level $j\geq1$ it chooses one of
$d_j$ labels uniformly and an independent uniform sign.  Its value is
then the corresponding unitary.  This description gives law $\mu$ even
when distinct labels happen to represent the same element of $M$.

Sample $\ell$ independent paths of length $n_\ell$, involving
$K_\ell$ increments in all.  Let $\mathcal B_\ell$ be the union of
the following four events: a level greater than $\ell$ appears; a
level-zero increment lies outside $F_\ell$; a level-$\ell$ label is
repeated; or some path has no level-$\ell$ increment.  Union bounds give
\begin{align}
 \Pr(\mathcal B_\ell)
 &\leq K_\ell\sum_{j>\ell}p_j
       +K_\ell p_0\mu_0(\U(M)\setminus F_\ell)\notag\\
 &\quad+\frac{\binom{K_\ell}{2}}{d_\ell}
       +\ell(1-p_\ell)^{n_\ell}.
 \label{eq:walk-bad-events}
\end{align}
For the collision term, any specified pair of increments has
probability $p_\ell^2/d_\ell\leq1/d_\ell$ of using the same
level-$\ell$ label; the signs do not matter.  For the last term, a
specified path misses that level with probability
$(1-p_\ell)^{n_\ell}$.

Every term on the right of \eqref{eq:walk-bad-events} tends to zero.
Indeed, $p_{\ell+1}=p_\ell^2$ and $p_\ell\leq1/4$, so
\[
 \sum_{j>\ell}p_j
 \leq\frac{p_\ell^2}{1-p_\ell^2}
 \leq\frac{16}{15}p_\ell^2.
\]
Consequently, the four terms admit the respective bounds
\begin{align*}
 \frac{16}{15}\ell\sqrt{p_\ell},\qquad
 K_\ell^{-1},\qquad
 \frac{1}{2K_\ell^2},\qquad
 \ell\exp(-p_\ell^{-1/2}),
\end{align*}
all converging to zero.

Outside $\mathcal B_\ell$, the sampled paths form exactly one of the
lists tested at stage $\ell$.  All their earlier-level factors belong
to the coefficient alphabet, each path has a level-$\ell$ label, and
all occurrences of those labels are distinct.  Thus, simultaneously
for every $w\in\mathbb F_\ell$ with $1\leq|w|\leq\ell$, the trace
of the evaluated word has absolute value less than $1/\ell$.

Now fix $r$, a nonempty finite $\mathcal W$, and $\eta>0$.  For
$\ell\geq\max\{r,\max_{w\in\mathcal W}|w|\}$, append independent
paths to the given $r$ paths to obtain $\ell$ paths.  On the preceding
good event every word in $\mathcal W$ has absolute trace below
$1/\ell$.  Once $1/\ell<\eta$, the probability in the theorem is at
most $\Pr(\mathcal B_\ell)$ and therefore tends to zero.

The measure $\mu$ is symmetric and countably supported by construction.
Its base component has positive weight, so each $v_j$ is a
positive-mass atom.  Since bounded $2$-norm convergence implies strong
convergence in the tracial representation, the von Neumann algebra
generated by these atoms contains every unitary and hence is $M$.
Finally, the identity group word always has trace one, and every
$*$-monomial in unitaries reduces to a group word.  The asserted
convergence therefore describes all joint $*$-moments.
\end{proof}

\begin{remark}\label{rem:walk-only-moments}
Theorem~\ref{thm:walk} asserts convergence of scalar moments.  It does
not assert convergence of operator norms of polynomials in the sampled
products.  Its finite-family formulation permits further conditions
whose failure probabilities tend to zero to be imposed simultaneously
at any fixed stage of a later selection argument.
\end{remark}
